\section{问题三的建模与求解}

本节先明确门店布局重要性的业务含义和评价对象，再构造指标与权重，利用综合评价模型排序，最后检验排名稳定性。

\subsection{问题分析}
说明评价对象、渠道分组、时间和区域维度，以及“重要性”的业务定义。

\subsection{评价对象与指标体系}
说明指标来源、正负方向、重复性检查和业务含义。

\subsection{标准化、赋权与综合评价模型}
\subsubsection{指标标准化与权重}
说明标准化、权重来源和权重组合方式。

\subsubsection{综合评价模型}
给出评价公式、距离或得分定义。

\subsection{排序结果与推荐对象}
呈现分渠道排名、前三门店和结果解释。

\subsection{稳定性检验与运营建议}
说明权重扰动、替代口径和排名变化，并给出简短建议。

\subsection{问题三小结}
直接列出线上、线下推荐对象及稳定性结论。
